\subsection{A global version of the Springer action}

The product formula in Theorem \ref{th:prod} suggests that there should be a global analog of affine Springer theory where affine Springer fibers are replaced by Hitchin fibers. Such a theory was developed in a series of papers of the author starting with \cite{GS}.

\subsubsection{Iwahori level structure} We now define a Hitchin-type analog of the affine Springer fibers $\sY_{\gamma}$. We fix the curve $X$ and a line bundle $\calL$ on it as before. Let $\Mpar_{G,\calL}$ be the moduli stack classifying $(x,\calE,\varphi, \calE^{B}_{x})$ where $x\in X$, $(\calE,\varphi)$ is an $\calL$-twisted Higgs $G$-bundle, and $\calE^{B}_{x}$ is a reduction of the fiber $\calE_{x}$ at $x$ to a $B$-torsor $\calE^{B}_{x}$ (here $B\subset G$ is a Borel subgroup of $G$) compatible with $\varphi$. 

When $G=\GL_{n}$, $\Mpar_{G,\calL}$ classifies, in addition to an $\calL$-twisted Higgs bundle $(\calV,\varphi)$, a point $x\in X$ and a full flag of the fiber $\calV_{x}$. Such a full flag is the same data as a chain of coherent subsheaves $\calV_{i}\subset\calV$
\begin{equation*}
\calV(-x)=\calV_{0}\subset \calV_{1}\subset\cdots\subset\calV_{n-1}\subset\calV_{n}=\calV
\end{equation*}
such that $\calV_{i}/\calV_{i-1}$ has length $1$ supported at $x$. The compatibility condition between $\varphi$ and the full flag requires that $\varphi$ restrict to a map $\calV_{i}\to\calV_{i}\otimes\calL$ for each $0\leq i\leq n$. 

We have an analog of the Hitchin fibration
\begin{equation*}
\fpar_{G,\calL}:\Mpar_{G,\calL}\to \calA_{G,\calL}\times X
\end{equation*}
that records also the point $x\in X$ in the data, in addition to $a=f_{G,\calL}(\calE,\varphi)\in\calA_{G,\calL}$.

\begin{theorem}[See {\cite[Th 3.3.3]{GS}}]\label{th:GS} Suppose $\deg\calL>2g-1$, then there is a natural action of $\tilW=\xcoch(T)\rtimes W$ on the restriction of the complex $\bR\fpar_{!}\Ql$ to a certain open subset $(\Ah_{G,\calL}\times X)'$ of $\Ah_{G,\calL}\times X$. \footnote{When $\textup{char}(k)=0$, one can take $(\Ah_{G,\calL}\times X)'=\Ah_{G,\calL}\times X$; in general, the restriction to the open subset $(\Ah_{G,\calL}\times X)'$ does not limit applications to questions about affine Springer fibers.}
\end{theorem}

\subsubsection{Extended symmetry} Just as in the case of affine Springer fibers, the $\tilW$-action in Theorem \ref{th:GS} may be extended to an action of the graded version of the double affine Hecke algebra, see \cite[Th 6.1.6]{GS}. Also, there is a product formula relating the fiber $\Mpar_{a,x}$ over $(a,x)\in\Ah_{G,\calL}\times X$ and the product of the affine Springer fiber $\sY_{a_{x}}$ with affine Springer fibers $\sX_{a_{y}}$ for $y\neq x$.  The induced $\tilW$-action on the stalk $(\bR\fpar_{!}\Ql)_{a,x}\cong\cohog{*}{\Mpar_{a,x}}$ is compatible with the affine Springer action on $\cohog{*}{\sY_{a_{x}}}$ in Theorem \ref{th:aff action}. This connection between $\tilW$-actions on the cohomology of Hitchin fibers and affine Springer fibers can be used to prove results about affine Springer actions. See \cite{YSph} for such an application.




\subsection{Exercises}

\subsubsection{} Suppose $G=\SL_{n}$. Describe the Hitchin fibers over $\Ah_{G,\calL}$ in terms of spectral curves.

\subsubsection{} Suppose $G=\SL_{n}$. Compute the dimension of $\calA_{G,\calL}$ and of a Hitchin fiber $\calM_{a}$. When $\calL=\omega_{X}$, check that they have the same dimension. This is a numerical evidence that the Hitchin fibration in this case is a Lagrangian fibration.


\subsubsection{} Describe the Hitchin base for $G=\SO_{n}$.

\subsubsection{} For $G=\SO_{n}$, describe Hitchin fibers in terms of spectral curves.

\subsubsection{} Let $C$ be a rational curve over $k$ with a unique singularity $x_{0}$ which is unibranched (i.e., the preimage of $x_{0}$ in the normalization $\tilC\cong\PP^{1}$ is a single point).  Let $\overline{P}_{x_{0}}$ be the moduli space of fractional ideals for the completed local ring $\wh\calO_{C,x_{0}}$. Show that there is a canonical homeomorphism
\begin{equation*}
\overline{P}_{x_{0}}\to \cPic(C).
\end{equation*}
Explain why this is a special case of Proposition \ref{p:Pic prod}.

